Data centers are the material substrate of the digital economy — responsible for an estimated 400~TWh of electricity consumption annually and constituted by dense constellations of corporate, financial, civic, and governmental actors. Yet a bibliometric audit of 4,757 publications on data centers (OpenAlex, 2010--2025) reveals a striking blind spot: not one engages the Multi-Level Perspective, Technological Innovation Systems, Strategic Niche Management, or sociotechnical regimes frameworks that are standard equipment in innovation and sustainability transitions research. The literature fragments into three non-communicating silos — technical efficiency, environmental impact, and socio-political dynamics — none of which can account for the systemic, multi-actor character of the phenomena they study.

This paper proposes \textit{datacentering} as a conceptual reframing — a gerund that foregrounds the active, contested, ongoing process of assembling digital infrastructure rather than treating it as a stable technical object. Drawing on complex adaptive systems (CAS) thinking as a sensemaking lens, we develop a multimodal research agenda structured around the conceptual and structural dimensions of data center assemblages: how actors construct the imaginaries that make infrastructure legible or illegible, and how ownership hierarchies and actor constellations shape governance capacity. We propose four complementary methodological work packages — corporate and territorial text analysis, visual and satellite imagery analysis, ownership cartography, and contestation signal detection — and illustrate the analytical purchase of the framework through a comparison of Northern Virginia and Amsterdam.

We conclude that the complexity of data center assemblages as CAS-like phenomena demands not only multimodal methods but a collective, open approach to knowledge production. The paper is accompanied by an open research infrastructure — a community-maintained facility index, shared analytical pipelines, and open data — on the grounds that no single researcher or team can build this cartography alone.
